% Options for packages loaded elsewhere
\PassOptionsToPackage{unicode}{hyperref}
\PassOptionsToPackage{hyphens}{url}
\PassOptionsToPackage{dvipsnames,svgnames,x11names}{xcolor}
%
\documentclass[
  11pt,
]{article}
\usepackage{amsmath,amssymb}
\usepackage{iftex}
\ifPDFTeX
  \usepackage[T1]{fontenc}
  \usepackage[utf8]{inputenc}
  \usepackage{textcomp} % provide euro and other symbols
\else % if luatex or xetex
  \usepackage{unicode-math} % this also loads fontspec
  \defaultfontfeatures{Scale=MatchLowercase}
  \defaultfontfeatures[\rmfamily]{Ligatures=TeX,Scale=1}
\fi
\usepackage{lmodern}
\ifPDFTeX\else
  % xetex/luatex font selection
\fi
% Use upquote if available, for straight quotes in verbatim environments
\IfFileExists{upquote.sty}{\usepackage{upquote}}{}
\IfFileExists{microtype.sty}{% use microtype if available
  \usepackage[]{microtype}
  \UseMicrotypeSet[protrusion]{basicmath} % disable protrusion for tt fonts
}{}
\makeatletter
\@ifundefined{KOMAClassName}{% if non-KOMA class
  \IfFileExists{parskip.sty}{%
    \usepackage{parskip}
  }{% else
    \setlength{\parindent}{0pt}
    \setlength{\parskip}{6pt plus 2pt minus 1pt}}
}{% if KOMA class
  \KOMAoptions{parskip=half}}
\makeatother
\usepackage{xcolor}
\usepackage[margin=1in]{geometry}
\usepackage{longtable,booktabs,array}
\usepackage{calc} % for calculating minipage widths
% Correct order of tables after \paragraph or \subparagraph
\usepackage{etoolbox}
\makeatletter
\patchcmd\longtable{\par}{\if@noskipsec\mbox{}\fi\par}{}{}
\makeatother
% Allow footnotes in longtable head/foot
\IfFileExists{footnotehyper.sty}{\usepackage{footnotehyper}}{\usepackage{footnote}}
\makesavenoteenv{longtable}
\setlength{\emergencystretch}{3em} % prevent overfull lines
\providecommand{\tightlist}{%
  \setlength{\itemsep}{0pt}\setlength{\parskip}{0pt}}
\setcounter{secnumdepth}{-\maxdimen} % remove section numbering
\ifLuaTeX
  \usepackage{selnolig}  % disable illegal ligatures
\fi
\IfFileExists{bookmark.sty}{\usepackage{bookmark}}{\usepackage{hyperref}}
\IfFileExists{xurl.sty}{\usepackage{xurl}}{} % add URL line breaks if available
\urlstyle{same}
\hypersetup{
  pdfauthor={Stefano Alimonti --- stefalimonti@gmail.com},
  colorlinks=true,
  linkcolor={blue},
  filecolor={Maroon},
  citecolor={Blue},
  urlcolor={blue},
  pdfcreator={LaTeX via pandoc}}

\author{Stefano Alimonti --- stefalimonti@gmail.com}
\date{March 2026}

\begin{document}

{
\hypersetup{linkcolor=black}
\setcounter{tocdepth}{2}
\tableofcontents
}
\section{Reader's Guide: How π Emerges from Binary
Multiplication}\label{readers-guide-how-ux3c0-emerges-from-binary-multiplication}

\textbf{Stefano Alimonti} --- March 2026

\emph{This guide explains, step by step, how the number π = 3.14159\ldots{}
appears in a problem that has nothing to do with circles. It requires no
mathematical background beyond knowing how to multiply.}

\begin{center}\rule{0.5\linewidth}{0.5pt}\end{center}

\newpage

\subsection{1. Multiplying in Binary}\label{multiplying-in-binary}

Computers multiply numbers in binary (base 2), using only the digits 0 and 1.
The same rules apply as in decimal --- multiply each digit, add the columns,
carry the overflow.

Here is \(13 \times 11 = 143\) in binary:

\begin{verbatim}
          1 1 0 1          (13 in binary)
        × 1 0 1 1          (11 in binary)
        ---------
          1 1 0 1          13 × 1
        1 1 0 1 ·          13 × 1, shifted left
      0 0 0 0 · ·          13 × 0, shifted left twice
    1 1 0 1 · · ·          13 × 1, shifted left three times
    -----------------
    1 0 0 0 1 1 1 1        (143 in binary)
\end{verbatim}

Each column is added up. When a column totals more than 1, the excess is
\emph{carried} to the next column --- exactly like carrying in decimal addition.

\begin{center}\rule{0.5\linewidth}{0.5pt}\end{center}

\newpage

\subsection{2. The Carry Chain}\label{the-carry-chain}

Reading the columns from right to left, the column sums and carries for
\(13 \times 11\) are:

\begin{verbatim}
Column:      0   1   2   3   4   5   6
Sum:         1   1   1   3   1   1   1
Carry in:    0   0   0   0   1   1   1
Total:       1   1   1   3   2   2   2
Output bit:  1   1   1   1   0   0   0
Carry out:   0   0   0   1   1   1   1
\end{verbatim}

The \emph{carry chain} --- the sequence 0, 0, 0, 0, 1, 1, 1, 1 --- records the
overflow at each step. At column 3, the sum exceeds 1 for the first time, and a
carry is born. Once present, it persists.

The key insight from probability theory (Diaconis and Fulman, 2009): if the
input digits are random, the carry at each position depends only on the
\emph{current} carry and the column sum --- not on anything earlier. This makes
the carry chain a \textbf{Markov process}, amenable to spectral analysis.

\begin{center}\rule{0.5\linewidth}{0.5pt}\end{center}

\newpage

\subsection{3. How Many Digits Does the Product
Have?}\label{how-many-digits-does-the-product-have}

When we multiply two 4-bit numbers (8 through 15), the product can have either 7
or 8 bits:

\begin{itemize}
\tightlist
\item
  \(8 \times 8 = 64\) → 7 bits (1000000)
\item
  \(11 \times 11 = 121\) → 7 bits (1111001)
\item
  \(12 \times 12 = 144\) → 8 bits (10010000)
\item
  \(15 \times 15 = 225\) → 8 bits (11100001)
\end{itemize}

The boundary is at 128: products below 128 have 7 bits, products at or above 128
have 8 bits.

We call the 7-bit products \textbf{D-odd} (the product has an \emph{odd} number
of digits, \(D = 2K - 1 = 7\)). For D-odd products, the carry chain starts and
ends at zero --- a ``bridge'' that goes up and comes back down, like a ball
thrown into the air.

\begin{verbatim}
Carry value
    2 ┤
    1 ┤          ╭───╮
    0 ┤──────╮   │   ╰──── 0
             ╰───╯
      0  1  2  3  4  5  6  position
         D-odd carry bridge
\end{verbatim}

\begin{center}\rule{0.5\linewidth}{0.5pt}\end{center}

\newpage

\subsection{4. The Four Sectors}\label{the-four-sectors}

Each 4-bit number has the form \(1abc\) in binary. The \textbf{second-highest
bit} \(a\) divides the numbers into two groups:

\begin{itemize}
\tightlist
\item
  \textbf{Low group} (\(a = 0\)): 8, 9, 10, 11 (binary 1000 through 1011)
\item
  \textbf{High group} (\(a = 1\)): 12, 13, 14, 15 (binary 1100 through 1111)
\end{itemize}

Since we have two numbers \(X\) and \(Y\), their second bits define four
\textbf{sectors}:

\begin{verbatim}
            Y: low (8–11)       Y: high (12–15)
           ┌────────────────┬────────────────┐
X: low     │                │                │
(8–11)     │  Sector (0,0)  │  Sector (0,1)  │
           │  16 pairs      │  16 pairs      │
           │  16 D-odd ✓    │  8 D-odd  ✓    │
           ├────────────────┼────────────────┤
X: high    │                │                │
(12–15)    │  Sector (1,0)  │  Sector (1,1)  │
           │  16 pairs      │  16 pairs      │
           │  8 D-odd  ✓    │  0 D-odd  ✗    │
           └────────────────┴────────────────┘
\end{verbatim}

\textbf{A proven fact (Theorem 4):} Sector (1,1) has \emph{zero} D-odd pairs.
When both numbers are in the high group, the product always overflows to 8 bits.
This is not a coincidence --- it holds for all \(K\), because if
\(X \geq 3 \cdot 2^{K-2}\) and \(Y \geq 3 \cdot 2^{K-2}\) (both second bits are
1), then \(XY \geq 9 \cdot 2^{2K-4} > 2^{2K-1}\), so the product always has
\(2K\) bits.

\begin{center}\rule{0.5\linewidth}{0.5pt}\end{center}

\newpage

\subsection{5. The Carry Weight and the Sector
Ratio}\label{the-carry-weight-and-the-sector-ratio}

For each D-odd pair, we read the carry chain and assign a \textbf{weight} of
\(+1\), \(0\), or \(-1\). The ``cascade'' method scans the carry chain from top
to bottom and reads the carry value \(c\) just below the highest nonzero carry
position. The weight is \(c - 1\): if \(c = 0\), the weight is \(-1\); if
\(c = 1\), the weight is \(0\); if \(c \geq 2\), the weight is positive.

We add up these weights across all D-odd pairs in each sector: - \(\sigma_{00}\)
= total weight from sector (0,0) - \(\sigma_{10}\) = total weight from sector
(1,0)

The \textbf{sector ratio} is their quotient:

\[R(K) = \frac{\sigma_{10}}{\sigma_{00}}\]

This is a purely combinatorial quantity: it counts carry patterns in binary
multiplication, weighted by the cascade valuation. No geometry, no circles, no
angles.

\begin{center}\rule{0.5\linewidth}{0.5pt}\end{center}

\newpage

\subsection{6. The Punchline: R(K) Converges to
−π}\label{the-punchline-rk-converges-to-ux3c0}

We computed \(R(K)\) for every value of \(K\) from 4 to 21 by exhaustive
enumeration --- checking every single pair of \(K\)-bit numbers. At \(K = 21\),
this means examining over \textbf{one trillion} pairs.

\begin{verbatim}
R(K)
 +2/3 ─ ─ ─ ─ ─ ─ ─ ─ ─ ─ ─ ─ ─  ← Markov prediction (rational)
  0   ┤
 -1   ┤     ·
 -2   ┤          ·
 -3   ┤               · · · · · ·  ← converging
−π  ─ ┤─ ─ ─ ─ ─ ─ ─ ─ ─ ─ ─ ─ ·  ← limit = −3.14159...
 -4   ┤
      └──┬──┬──┬──┬──┬──┬──┬──┬──
         7  9 11 13 15 17 19 21   K
\end{verbatim}

\begin{longtable}[]{@{}llll@{}}
\toprule\noalign{}
\(K\) & Pairs checked & \(R(K)\) & Gap from \(-\pi\) \\
\midrule\noalign{}
\endhead
\bottomrule\noalign{}
\endlastfoot
7 & 4,096 & \(-0.09\) & 3.05 \\
11 & \textasciitilde1 million & \(-1.55\) & 1.59 \\
15 & \textasciitilde270 million & \(-2.82\) & 0.32 \\
19 & \textasciitilde69 billion & \(-3.110\) & 0.032 \\
21 & \textasciitilde1.1 trillion & \(-3.133\) & 0.008 \\
\(\infty\) (extrapolated) & & \(-3.1416\ldots\) & \(< 0.001\) \\
\end{longtable}

After Richardson extrapolation (a standard acceleration technique), the limit
matches \(-\pi = -3.14159\ldots\) to \textbf{4 significant digits}.

\begin{center}\rule{0.5\linewidth}{0.5pt}\end{center}

\newpage

\subsection{7. Why Is This Surprising?}\label{why-is-this-surprising}

Three things make this result remarkable:

\textbf{1. The Markov model is completely wrong.} If carries at each position
were independent (the Markov assumption), the sector ratio would approach
\(R = +2/3\) --- a positive rational number. The observed answer instead points
toward \(-\pi\) --- negative, irrational, transcendental. The sign is wrong, the
magnitude is wrong by a factor of 5, and the number itself is of a fundamentally
different mathematical nature.

\textbf{2. π has no reason to be here.} The number π is the ratio of a circle's
circumference to its diameter. Our problem involves no circles, no angles, no
curves --- just binary digits, column sums, and carries. The connection to π is
deep: it enters through the Dirichlet character \(\chi_4\), which satisfies
\(L(1, \chi_4) = \pi/4\) (the Leibniz series
\(1 - 1/3 + 1/5 - 1/7 + \cdots = \pi/4\)). The carry chain's spectral structure
naturally selects this character.

\textbf{3. A sharp phase transition is proposed.} In the LMH model of {[}E{]},
increasing the correlation between adjacent carries continuously transforms the
spectral sum from \(+2/3\) (Markov, subcritical) through \(0\) (critical point)
to \(-\pi\) (supercritical). The governing Möbius transformation is

\[S(A) = \frac{2(1-A)}{3-2A}\]

where \(A\) measures the strength of digit correlations. Matching this model sum
to \(-\pi\) gives \(A^* = (3\pi + 2)/(2(\pi + 1)) \approx 1.379\). If the
physical sector ratio is governed by the LMH, this is the corresponding critical
value.

\begin{center}\rule{0.5\linewidth}{0.5pt}\end{center}

\newpage

\subsection{8. What Is Proved and What Is
Conjectured}\label{what-is-proved-and-what-is-conjectured}

\begin{longtable}[]{@{}
  >{\raggedright\arraybackslash}p{(\columnwidth - 2\tabcolsep) * \real{0.5000}}
  >{\raggedright\arraybackslash}p{(\columnwidth - 2\tabcolsep) * \real{0.5000}}@{}}
\toprule\noalign{}
\begin{minipage}[b]{\linewidth}\raggedright
Result
\end{minipage} & \begin{minipage}[b]{\linewidth}\raggedright
Status
\end{minipage} \\
\midrule\noalign{}
\endhead
\bottomrule\noalign{}
\endlastfoot
The closed-form Möbius transformation \(S(A) = 2(1-A)/(3-2A)\) & \textbf{Proved}
(Theorem 1) \\
The Markov baseline \(R_{\text{Markov}}(L) = +2/3 + O(L^{-2})\) &
\textbf{Proved} (Theorem 2) \\
The spectral zeta function reproduces \(\zeta(2s)\) & \textbf{Proved} (Theorem
3) \\
No D-odd pair has sector (1,1) & \textbf{Proved} (Theorem 4) \\
The count ratio \(N_{10}/N_{00}\) is a closed-form logarithmic constant &
\textbf{Proved} (Theorem 5) \\
Cascade stopping depths in sector (0,0) obey structural rigidity &
\textbf{Proved} (Theorem 6) \\
The first nonzero carry is exactly 1 & \textbf{Proved} (Theorem 7) \\
\(R(\infty) = -\pi\) & \textbf{Conjectured} (4.0-digit numerical evidence) \\
\end{longtable}

The conjecture reduces to proving two things: (i) a scalar constant equals
\(-4\), and (ii) the spectral gap \(1/2\) is preserved under the D-odd boundary
condition. The companion paper {[}L{]} now realizes the same stopping-time
channel as a canonical weighted first-return resolvent and shows that the
current carry-side analytic object captures Euler/local-factor structure but
does not yet transfer the \(L\)-function zeros. The algebraic framework is
complete; the remaining step is analytical.

\begin{center}\rule{0.5\linewidth}{0.5pt}\end{center}

\newpage

\subsection{9. The Bigger Picture}\label{the-bigger-picture}

The conjectural limit \(R \to -\pi\) is one facet of a broader connection
between binary carries and number theory:

\begin{itemize}
\tightlist
\item
  The \textbf{Riemann zeta function} \(\zeta(s)\), which encodes the
  distribution of primes, is approximated by ensemble-averaged spectral
  determinants of carry companion matrices {[}B{]}.
\item
  The \textbf{spectral zeta function} of the Markov carry bridge reproduces
  \(\zeta(2s)\) at all even arguments (Theorem 3 of this paper).
\item
  Conjecturally, the sector ratio converges to \(-4L(1, \chi_4) = -\pi\),
  connecting to the \textbf{Dirichlet \(L\)-function} of the character modulo 4.
\item
  Together, \(\zeta(s)\) and \(L(s, \chi_4)\) form the \textbf{Dedekind zeta
  function} of the Gaussian integers \(\mathbb{Z}[i]\):
  \(\zeta_{\mathbb{Q}(i)}(s) = \zeta(s) \cdot L(s, \chi_4)\).
\end{itemize}

Binary multiplication, through its carry structure, already produces one factor
of this Dedekind zeta function via the carry polynomial ensemble {[}B{]}, and
conjecturally contributes the second through the sector ratio studied here.

\begin{center}\rule{0.5\linewidth}{0.5pt}\end{center}

\newpage

\subsection{Glossary}\label{glossary}

\begin{itemize}
\tightlist
\item
  \textbf{Binary (base 2):} A number system using only digits 0 and 1.
\item
  \textbf{Carry:} The overflow when a column sum exceeds the base. In base 2, a
  column sum of 2 produces a carry of 1.
\item
  \textbf{D-odd:} A product whose binary representation has exactly \(2K - 1\)
  digits (an odd number).
\item
  \textbf{Sector:} One of four groups determined by the second-highest bit of
  each factor.
\item
  \textbf{Sector ratio \(R(K)\):} The ratio of carry-weighted sums in sectors
  (1,0) and (0,0).
\item
  \textbf{Markov process:} A random process where the next state depends only on
  the current state.
\item
  \textbf{Phase transition:} A qualitative change in behavior as a parameter
  crosses a threshold.
\item
  \textbf{Richardson extrapolation:} A numerical technique that accelerates
  convergence by eliminating leading error terms.
\end{itemize}

\begin{center}\rule{0.5\linewidth}{0.5pt}\end{center}

\newpage

\subsection{References}\label{references}

\begin{enumerate}
\def\labelenumi{\arabic{enumi}.}
\tightlist
\item
  P. Diaconis, J. Fulman, ``Carries, Shuffling, and Symmetric Functions,''
  \emph{Adv. Appl. Math.} 43(2), 176--196, 2009.
\item
  {[}A{]} S. Alimonti, ``Spectral Theory of Carries in Positional
  Multiplication,'' this series.
  doi:\href{https://doi.org/10.5281/zenodo.18895593}{10.5281/zenodo.18895593}
  ---
  \href{https://github.com/stefanoalimonti/carry-arithmetic-A-spectral-theory}{GitHub}
\item
  {[}B{]} S. Alimonti, ``Carry Polynomials and the Euler Product,'' this series.
  doi:\href{https://doi.org/10.5281/zenodo.18895597}{10.5281/zenodo.18895597}
  ---
  \href{https://github.com/stefanoalimonti/carry-arithmetic-B-zeta-approximation}{GitHub}
\item
  {[}E{]} S. Alimonti, ``The Trace Anomaly of Binary Multiplication,'' this
  series.
  doi:\href{https://doi.org/10.5281/zenodo.18895604}{10.5281/zenodo.18895604}
  ---
  \href{https://github.com/stefanoalimonti/carry-arithmetic-E-trace-anomaly}{GitHub}
\end{enumerate}

\begin{center}\rule{0.5\linewidth}{0.5pt}\end{center}

\emph{CC BY 4.0}

\end{document}
